(1) 局地垂直自转分量为 $\omega_F=\Omega_\oplus\sin\lambda$，摆的小振动频率为 $\omega_0=\sqrt{g/L}$。水平科里奥利加速度 $-2\boldsymbol\Omega\times\mathbf v$ 给 $\ddot x+\omega_0^2x=2\omega_F\dot y$、$\ddot y+\omega_0^2y=-2\omega_F\dot x$；另两个水平自转分量只通过已忽略的竖直耦合进入。
(2) 令 $z=x+iy$，得 $\ddot z+2i\omega_F\dot z+\omega_0^2z=0$。变换 $z=e^{-i\omega_F t}w$ 后，$\ddot w+(\omega_0^2+\omega_F^2)w=0$，故 $z=e^{-i\omega_Ft}[C\cos(\omega_bt)+D\sin(\omega_bt)]$，其中 $\omega_b=\sqrt{\omega_0^2+\omega_F^2}$，$C,D$ 由复初值给定。
(3) 若在慢旋坐标下取线性极化振荡，$w=A\cos(\omega_bt)$，则 $z=Ae^{-i\omega_Ft}\cos(\omega_bt)$，摆线方向角为 $-\omega_Ft$（模 $\pi$）。因此北半球摆动平面相对地面顺时针缓慢进动，大小 $|\Omega_\oplus\sin\lambda|$；并非摆球的高频振动频率。
(4) 赤道 $\lambda=0$ 时不发生理想傅科进动，北极 $\lambda=\pi/2$ 时角速率为 $-\Omega_\oplus$，南半球 $\sin\lambda<0$ 时转向反号。推导依赖小振幅和 $\Omega_\oplus\ll\sqrt{g/L}$，大振幅的非线性摆频变化与竖向运动必须重新纳入。
